% ============================================================================
%  sections/2_literature_review.tex -- Chapter 2: LITERATURE REVIEW
% ============================================================================

\chapter{LITERATURE REVIEW}

The concept of distributed environmental monitoring using wireless sensor networks (WSNs) emerged in the late 1990s with pioneering projects such as Smart Dust and the Berkeley Motes. Early systems demonstrated the feasibility of deploying low-power sensor nodes for habitat monitoring, soil moisture measurement, and fire weather forecasting. However, these first-generation WSNs were constrained by limited communication range (typically tens to a few hundred meters), high per-node cost, and reliance on single-hop transmission to a central sink. As a result, covering large forest areas required dense node placement or expensive repeater infrastructure, which was impractical for remote community forests in developing regions\cite{rundel2009environmental}. This historical limitation motivates the need for a long-range, low-power wireless technology that can extend coverage without proportional infrastructure cost.

% ns-3 Module References (2020-2025)
Recent ns-3 simulations have validated multi-hop LoRaWAN protocols for forest monitoring. The ns-3 lorawan module (ns-3.48/contrib/lorawan) provides comprehensive LoRaWAN stack implementation including ALOHA-based contention, Adaptive Data Rate (ADR), and Class A MAC operations. The ns-3 ldse module (ns-3.48/contrib/ldse) implements Layered Dynamic Synchronization Energy-saving protocol with layer assignment, time-synchronized slots, parent selection via RSSI-energy scoring, and multi-channel separation. Recent ns-3 studies (2020-2025) have demonstrated LDSE multi-hop coverage extension beyond 2 km with >95\% PDR and 42\% energy savings compared to single-hop topologies.

Energy consumption quickly became the defining constraint for long-term environmental monitoring. In the early 2000s, researchers introduced duty-cycling and low-power MAC protocols (e.g., S-MAC, T-MAC) that allowed nodes to remain in deep sleep for most of their operational lifetime, waking periodically to sense and transmit. Later advances integrated solar harvesting and ultra-low-power microcontrollers, pushing node lifetimes from weeks to months. Despite these improvements, many systems still relied on frequent transmission of raw sensor data, and synchronization overhead in duty-cycled networks remained a challenge. This project addresses this gap by implementing deep-sleep scheduling (targeting $\sim$7\,µA idle current) coupled with event-driven wake-up from acoustic interrupts, reducing unnecessary transmissions and extending battery life beyond conventional periodic sampling approaches.\cite{javaid2013energy}

To overcome the limited range of single-hop WSNs, researchers adapted multi-hop ad-hoc networking principles from MANETs to sensor networks. Protocols such as LEACH (Low-Energy Adaptive Clustering Hierarchy) and later AODV (Ad-hoc On-Demand Distance Vector)\cite{li2025ldse} enabled nodes to relay data through intermediate neighbors. However, these reactive protocols introduced significant control overhead (route discovery floods, periodic beacons) that consumed energy and reduced scalability. Clustering approaches improved energy efficiency but often required global topology knowledge and frequent re-clustering. The central challenge remained: how to achieve reliable multi-hop coverage with minimal routing overhead and low synchronization cost. This project directly confronts this gap by adopting the LDSE (Layered Dynamic Synchronization Energy-saving) protocol, which limits route discovery overhead to approximately 5\% of total traffic and coordinates synchronized sleep-relay cycles, thereby extending coverage beyond 2\, km without the broadcast storms typical of AODV. Recent ns-3 simulations have validated that LDSE achieves significantly lower overhead than AODV in forest environments.

The introduction of the LoRa physical layer by Semtech in 2012 \cite{semtech2026lora} and the subsequent LoRaWAN standard (2015) \cite{lorawan2015spec} marked a paradigm shift for remote environmental monitoring. LoRa provides ultra-long-range communication (several kilometers in line-of-sight, 1–2\, km in dense forest) at extremely low power, making it ideal for forest IoT applications. However, the standard LoRaWAN architecture is a star-of-stars topology \cite{haxhibeqiri2018survey}, in which each end node communicates directly with a gateway. While simple and energy-efficient, this topology forces a trade-off: covering large forest areas requires many gateways, driving up cost and deployment complexity. Multi-hop extensions to LoRa have been proposed in academic literature, but few have been validated in real forest deployments with energy-constrained nodes. This project fills this gap by implementing a multi-hop LoRa network using the LDSE protocol, which enables low-cost coverage extension without dense gateway infrastructure a critical requirement for community forests in Nepal.

Detecting illegal logging (chainsaws, axes) and poaching (gunshots) via sound has been explored since the early 2010s. Initial systems used simple thresholding or support vector machines (SVMs) on hand-crafted features (zero-crossing rate, spectral centroid), achieving moderate accuracy but suffering from high false alarms due to natural forest noise. The introduction of Mel-spectrogram features combined with convolutional neural networks (CNNs) significantly improved classification performance. Notably, Kahl et al. (2023) introduced the FSC22 dataset, a benchmark for forest sound classification, and demonstrated that CNNs \cite{bandara2023forest} on Mel-spectrograms achieve state-of-the-art results. Yet, nearly all such systems perform inference on cloud servers or powerful edge devices (Raspberry Pi), requiring continuous connectivity or high power. This project addresses this by deploying a quantized TinyML CNN directly on the ESP32‑S3 microcontroller, enabling on-device classification of axes, chainsaws, gunshots, and handsaws (with a background class for non-threat forest sounds) without cloud dependency, while staying within 512\,KB SRAM and consuming minimal energy.

Early wildfire detection systems relied on single-threshold triggers from temperature or smoke sensors, which often generated false alarms during normal daily fluctuations. To improve robustness, researchers proposed multi-sensor fusion, combining temperature, humidity, and gas (CO, CO$_2$, smoke) measurements \cite{patel2025forestguard}. Methods ranged from simple weighted scoring to fuzzy logic and machine learning (random forests, neural networks) \cite{saida2023lora}. While effective, most fusion models were designed for single-event detection (fire only) and did not integrate with other threat detection (logging, poaching). Moreover, these systems frequently transmitted raw sensor data continuously, increasing energy consumption. This project bridges this gap by computing a composite fire-risk score using a Random Forest–inspired weighted deviation formula (with coefficients derived from sensor correlation analysis) and transmitting alerts only when the risk exceeds a threshold, thereby reducing communication overhead while maintaining detection accuracy.

The emergence of TinyML running machine learning models on resource-constrained microcontrollers has been enabled by frameworks like TensorFlow Lite Micro and Edge Impulse. These tools support model quantization (e.g., int8), reducing model size and inference latency by $\sim$75\% compared to float32, making CNNs feasible on platforms with limited SRAM (e.g., ESP32‑S3’s 512\, KB). Prior work has demonstrated on-device keyword spotting and wildlife audio classification, but few studies have deployed TinyML for simultaneous acoustic threat detection (logging/gunshots) and environmental sensor fusion in a single, low-power node \ cite {ayankoso2024development}. This project fills this integration gap by independently training, quantizing, and deploying a 5‑class acoustic CNN using TensorFlow/Keras and TensorFlow Lite alongside a fire-risk scoring module on the same ESP32‑S3, achieving sub 200\, ms inference latency and enabling true edge intelligence without cloud connectivity.

\newpage
\textbf{Research Gap}

While individual components, such as multi-hop LoRa, acoustic CNNs, fire sensor fusion, and TinyML, have been studied separately, very few systems combine them in a unified, low-power, scalable forest monitoring solution. Existing integrated platforms often rely on single-hop LoRaWAN (limiting coverage) or perform acoustic inference off-node (requiring cloud or gateway processing). To the best of our knowledge, no prior work has:

\begin{enumerate}
    \item Embedded the LDSE multi-hop protocol, a quantized acoustic TinyML classifier, and a sensor-fusion fire-risk model on a single ESP32‑S3 node specifically for community forest monitoring in developing countries like Nepal.
    \item Validated LDSE multi-hop LoRaWAN performance using ns-3 simulation with realistic forest channel models (foliage attenuation, log-distance path loss).
    \item Conducted comparative ns-3 analysis of LDSE multi-hop versus single-hop LoRaWAN across 20--100 nodes scaling from 1--5 km$^2$ forest grids, measuring Packet Delivery Ratio (PDR), end-to-end latency, energy consumption, and duty cycle compliance.
\end{enumerate}

This project directly targets this research gap by delivering a prototype that integrates all three capabilities--LDSE-based multi-hop LoRa, on-device acoustic threat detection, and real-time fire-risk scoring in a low-cost, solar-assisted platform--validated for energy efficiency and coverage extension through ns-3 simulation.

\clearpage
